\documentclass[11pt,a4paper]{article}
\usepackage[utf8]{inputenc}
\usepackage[T1]{fontenc}
\usepackage[french]{babel}
\usepackage{amsmath,amssymb,amsthm,geometry,hyperref}
\geometry{margin=2.4cm}
\setlength{\emergencystretch}{3em}
\newtheorem{proposition}{Proposition}
\newtheorem{corollary}{Corollaire}
\newcommand{\id}{I}
\title{Transport FE par espace auxiliaire avec mémoire\\
Stabilité, couplage et ordre temporel}
\author{Étude GPU SLFE, fork keckj/rodin}
\date{16 novembre 2026}
\begin{document}
\maketitle

\section{Le défaut introduit par les retours répétés}
Soient $V$ et $W$ deux espaces de Hilbert discrets contenant une constante
de même norme. Le transfert $R:V\to W$ est une contraction et vérifie
$R1_V=1_W$, $R^*1_W=1_V$. Les adjoints sont définis par les matrices de masse.
Pour un opérateur auxiliaire $S:W\to W$ conservatif et contractant,
le transport incrémental est
\[
 T(S)=\id+R^*(S-\id)R .
\]
La conservation et la contraction de cet opérateur ont été établies dans
la note sur le transfert incrémental.
L'identité $T(\id)=\id$ élimine une erreur fixe à pas nul. Elle ne donne
pas une propriété de semi-groupe. En effet,
\[
 T(S_2S_1)-T(S_2)T(S_1)
 =R^*(S_2-\id)(\id-RR^*)(S_1-\id)R .
\]
Le terme de droite mesure l'information auxiliaire supprimée entre les pas.

Si $S(t)=\exp(tL)$ avec $L^*=-L$, alors, à espaces fixés,
\[
 T(S(t))=\exp(tA)-\frac{t^2}{2}D+O(t^3),\quad
 A=R^*LR,\quad
 D=R^*L^*(\id-RR^*)LR\geq0 .
\]
Le terme quadratique dissipatif n'est généralement pas nul. Une composition
symétrique de sous-transports incrémentaux ne suffit donc pas à établir
un ordre temporel deux à maillage fixé.

\section{Une réalisation avec état auxiliaire persistant}
L'état initial est $v_0\in V$, $w_0=Rv_0\in W$. Un sous-pas d'advection
est réalisé par
\[
 w^+=Sw,\qquad v^+=v+R^*(w^+-w).
\]
Après une suite de sous-pas, le champ observé est
\[
 v_n=v_0+R^*(S_{n-1}\cdots S_0-\id)Rv_0 .
\]
Il est contractant comme application de $v_0$ vers $v_n$. Cette propriété
ne demande aucune nouvelle projection FE à chaque pas. L'énergie du champ
observé peut augmenter entre deux observations, tout en restant inférieure
à l'énergie initiale. Un aller-retour exactement inversé rend l'état initial;
son erreur finale ne mesure alors que la projection du champ initial.
Les essais d'advection doivent inclure des temps intermédiaires.

Cette forme persistante conserve un degré de liberté supplémentaire.
Pour $r=w-Rv$, le sous-pas donne
\[
 r^+=r+(\id-RR^*)(S-\id)w .
\]
Imposer $r=0$ après chaque pas restitue les transferts incrémentaux répétés
et leur défaut de composition.

\section{Coupler une opération FE sans perdre la mémoire}
Soit $G:V\to V$ une opération physique FE, linéaire et contractante.
Le sous-pas suivant applique exactement $G$ au champ FE,
\[
 v^+=Gv,\qquad w^+=w+R(v^+-v).
\]
Ce sous-pas laisse $r=w-Rv$ inchangé.
Sa stabilité avec les sous-pas d'advection ne résulte pas d'une monotonie
de $\|v\|$ à chaque advection. Elle résulte d'une énergie dans un espace
augmenté, dont le calcul n'est pas requis par l'algorithme.

\begin{proposition}
Posons $D_0=(\id-R^*R)^{1/2}$ et
$Jv=(Rv,D_0v)\in W\oplus V$.
Alors $J^*J=\id$. Le schéma à deux variables ci-dessus est l'observation
$v=J^*Z$ de la composition
\[
 Z\longmapsto
 \begin{pmatrix} S&0\\0&\id\end{pmatrix}Z,
 \qquad
 Z\longmapsto \bigl[\id+J(G-\id)J^*\bigr]Z,
 \qquad Z_0=Jv_0 .
\]
Ces deux opérateurs sont contractants. Si $S$ et $G$ préservent la masse et
la constante, alors le schéma les préserve également et
$\|v_n\|_V\leq\|Z_n\|\leq\|v_0\|_V$.
\end{proposition}
\begin{proof}
L'identité d'isométrie découle de
$R^*R+D_0^2=\id$. Écrivons $Z=(w,z)$.
L'advection donne directement
$J^*Z^+=v+R^*(S-\id)w$.
Pour $\delta=Gv-v$, le second opérateur donne
$w^+=w+R\delta$, $z^+=z+D_0\delta$ et
$J^*Z^+=v+\delta$.
La décomposition orthogonale
$Z=Jv+Z_\perp$ donne
$Z^+=JGv+Z_\perp$; sa norme ne croît pas.
La première application est contractante par hypothèse sur $S$.
Enfin $D_01_V=0$. La masse augmentée, associée au vecteur
$J1_V=(1_W,0)$, coïncide avec la masse observée et est préservée.
La même relation assure la préservation de la constante à partir
de l'état compatible $J1_V$.
\end{proof}

La racine $D_0$ sert uniquement à la preuve. Le code
\texttt{memory\_transport.py} utilise $R$, $R^*$, $v$ et $w$.
Les racines et exponentielles de matrices denses dans
\texttt{memory\_study.py} sont des oracles de validation GPU.
Le cas d'une application non linéaire $G$ vérifiant
$\|Gv\|\leq\|v\|$ conserve la borne de norme par la même décomposition.
Une contraction des différences demande l'hypothèse correspondante sur $G$.

\section{Conséquence pour l'ordre temporel}
Si $G(t)=\exp(tL_V)$, l'isométrie de $J$ donne exactement
\[
 \id+J(G(t)-\id)J^*=\exp(tJL_VJ^*).
\]
L'advection persistante et la physique relevée possèdent donc de vrais
générateurs dans l'espace augmenté. Une composition de Strang y est
d'ordre deux lorsque les sous-flots sont exacts ou suffisamment précis.
Cette affirmation porte sur le problème augmenté à espaces fixés.
Sa consistance avec l'EDP physique reste une question spatiale distincte.
Pour une physique dissipative, des pas négatifs ne sont généralement pas
contractants. Les compositions d'ordre quatre à coefficients négatifs ne
sont utilisées ici que pour les sous-transports unitaires.

\section{Borne conditionnelle sans accumulation des échanges}
Soient $v(t)\in V$, $y(t)\in W$ des approximations spatiales d'une solution
physique, $v_0=v(0)$ et $\mathcal S_n=S_{n-1}\cdots S_0$ contractant.
Posons $e(t)=v(t)-R^*y(t)$ et $\delta_0=Rv_0-y(0)$.
Un calcul exact donne
\[
 v_n-v(t_n)=e(0)-e(t_n)
 +R^*(\mathcal S_n-\id)\delta_0
 +R^*(\mathcal S_n y(0)-y(t_n)).
\]
Par conséquent,
\[
 \|v_n-v(t_n)\|
 \leq\|e(0)\|+\|e(t_n)\|+2\|\delta_0\|
       +\|\mathcal S_n y(0)-y(t_n)\| .
\]
La borne sépare les défauts spatiaux d'échange et l'erreur du solveur
auxiliaire. Aucun facteur égal au nombre de pas ne multiplie les trois
premiers défauts. Une estimation de convergence complète demande encore
des bornes pour ces défauts, adaptées à la géométrie et à la régularité
du problème. Cette proposition concerne l'advection persistante pure;
elle ne constitue pas à elle seule une estimation du problème couplé.

\section{Consistance conditionnelle du système couplé}
Considérons un problème spatial de référence
$\dot v=(A+L_V)v$ dans $V$, et un générateur auxiliaire $L_W$.
Supposons $L_W$ et $L_V$ dissipatifs dans leurs normes respectives.
Le générateur augmenté est
$\mathcal L=\operatorname{diag}(L_W,0)+JL_VJ^*$, donc dissipatif.
Pour la solution exacte augmentée $Z(t)$ issue de $Jv(0)$,
la différence $E(t)=Z(t)-Jv(t)$ est forcée par
\[
 f(t)=
 \begin{pmatrix}
  (L_WR-RA)v(t)\\-D_0Av(t)
 \end{pmatrix}.
\]
L'estimation d'énergie donne
\[
 \|J^*Z(t)-v(t)\|
 \leq \int_0^t
 \left(\|(L_WR-RA)v(s)\|^2+\|D_0Av(s)\|^2\right)^{1/2}\,ds.
\]
La physique FE n'introduit aucun défaut de transfert dans ce résidu.
La consistance est réduite à un défaut de commutation du transport et à
la résolution des dérivées transportées par le support auxiliaire.
Dans la factorisation $R=P_WP_Q|_V$, l'identité de Pythagore donne
\[
 \|D_0a\|^2
 =\|a-P_Qa\|^2+\|P_Qa-P_WP_Qa\|^2,\qquad a\in V.
\]
Ainsi ce second défaut a l'ordre des erreurs de projection, sans perte
supplémentaire d'un demi-ordre due à la racine carrée.
Une preuve de convergence vers l'EDP demande encore un choix de $A$,
ses estimations FE et une borne uniforme du commutateur.
Les résidus évalués sur un champ analytique initial sont des diagnostics
de cette condition, et non une preuve uniforme sur une trajectoire.

\section{Réalisation FE--DG--Fourier}
Dans les expériences, $V$ est DG $P_p$ sur triangles, $Q$ est DG $Q_q$
sur carrés fixes et $W$ est un espace trigonométrique sur le tore.
Le couplage est
\[
 R=P_WP_Q|_V,\qquad R^*=P_VP_Q|_W .
\]
Les trois projections sont orthogonales dans $L^2$. La factorisation est
donc contractante et préserve les constantes. Les intégrales FE--DG sont
statiques. Les moments DG--Fourier se factorisent en FFT des indices de
cellules et intégrales 1D de polynômes contre des exponentielles.
Ces dernières sont calculées par quadrature de Gauss d'ordre 32,
comparée à l'ordre 48 dans les validations. Les garanties exactes portent
sur les projections exactes; leur réalisation numérique est contrôlée
par l'adjonction, les constantes et les valeurs singulières.

Le transport test utilise
$\psi(x,y)=\sin^2(\pi x)\sin^2(\pi y)/\pi$ et
$u=(\partial_y\psi,-\partial_x\psi)$.
Ses quatre termes élémentaires sont des cisaillements dans les directions
cartésiennes et diagonales. Sur une grille uniforme impaire,
les changements $(x,y)\mapsto(x\pm y,y)$ sont des permutations.
Les translations par ligne sont des multiplications par des phases de
module un dans l'espace de Fourier. Chaque sous-transport est donc unitaire,
préserve la constante et la moyenne discrète. La norme de grille est la
norme $L^2$ exacte du polynôme trigonométrique représenté.
Les permutations diagonales peuvent replier les hautes fréquences;
l'unitarité ne prouve pas l'absence d'erreur d'aliasing.

Une composition symétrique des quatre modes est d'ordre deux pour le
générateur discret. La composition classique de Yoshida
$S_2(a t)S_2((1-2a)t)S_2(a t)$, $a=(2-2^{1/3})^{-1}$,
donne l'ordre quatre \cite{yoshida}. Chaque facteur demeure unitaire.
Les bornes de stabilité sont donc indépendantes du déplacement par pas;
la précision impose néanmoins une résolution temporelle des commutateurs.
Le facteur temporel commun $\cos(\pi t)$ est intégré analytiquement
avant cette composition. Une dépendance temporelle arbitraire de chaque
mode demande un traitement distinct.

\section{Portée et antécédents}
\subsection*{Directions de Fourier entières}
La réalisation ne se limite pas aux directions cartésiennes et diagonales.
La décomposition du flot en modes trigonométriques exactement intégrables
est un principe établi \cite{quispel}; la construction ci-dessous décrit
sa réalisation sur le support auxiliaire discret.
Pour un vecteur d'onde $k=g(a,b)$, avec $(a,b)$ primitif, le mode de vitesse
\[
u_k=\beta(-b,a)\sin(2\pi g(ax+by)+\theta)
\]
devient un cisaillement dans les coordonnées
$\xi=ax+by$, $\eta=cx+dy$, où les entiers $c,d$ satisfont $ad-bc=1$.
En effet, $\dot\xi=0$ et
$\dot\eta=\beta\sin(2\pi g\xi+\theta)$.
La matrice entière unimodulaire induit une permutation bijective de toute
grille périodique. Sur une grille impaire, le cisaillement de Fourier est
réel et unitaire. Une vitesse moyenne constante est transportée par une
translation de Fourier. Les compositions précédentes s'étendent donc à
une somme finie de tels modes. Le coût croît avec le nombre de directions;
les permutations peuvent toujours replier les fréquences. Cette extension
n'est pas un solveur efficace déjà démontré pour une vitesse FE quelconque.

\subsection*{Obstruction de résolution du support auxiliaire}
En advection pure, $v(t)-v_0=R^*(S(t)-I)Rv_0$. Par contraction,
\[
\|v(t)-v_0\|\leq 2\|Rv_0\|.
\]
Pour toute cible FE $\widehat v(t)$, l'inégalité triangulaire donne
\[
\|v(t)-\widehat v(t)\|\geq
\max\{0,\|\widehat v(t)-v_0\|-2\|Rv_0\|\}.
\]
Cette borne vaut pour tout nombre de pas et toute composition contractante.
Elle expose une limite structurelle : un champ bien représenté par le
maillage FE mais presque invisible sur le support auxiliaire ne peut pas
être correctement transporté. La stabilité indépendante du pas ne dispense
donc jamais d'une résolution spatiale auxiliaire adaptée.

Le transfert d'incréments, les dilatations d'opérateurs contractants,
les projections orthogonales et les compositions d'ordre élevé sont
des outils établis. Les transformations DG--Fourier sur GPU ont également
des antécédents explicites \cite{crews}. La décomposition en sous-champs
à divergence nulle appartient à la littérature des intégrateurs
préservant le volume \cite{xue}.
LEoPart réinjecte déjà des incréments de champs FE dans les propriétés
particulaires suivant une approche FLIP \cite{leopart}. Le retour d'un
incrément de physique FE vers un état auxiliaire n'est donc pas en soi
une nouveauté. La stabilité de cette réalisation précise et les estimations
ci-dessus doivent être distinguées de cette idée générale.
La nouveauté éventuelle doit être évaluée sur la réalisation FE avec
mémoire, sa consistance, ses couplages et ses coûts, et non sur ces outils
pris isolément.

Les expériences couvrent des triangles affines périodiques et une vitesse
prescrite à quatre modes. Les frontières non périodiques, les vitesses FE
quelconques, la compression, les garanties locales de flux, l'adaptation
du support, les couplages non linéaires et le calcul distribué ne sont pas
validés par ces résultats.

\begin{thebibliography}{9}
\bibitem{quispel} G. R. W. Quispel et D. I. McLaren,
\emph{Explicit volume-preserving and symplectic integrators for
trigonometric polynomial flows}, Journal of Computational Physics 186
(2003), 308--316,
\url{https://doi.org/10.1016/S0021-9991(03)00068-8}.
\bibitem{leopart} J. M. Maljaars, C. N. Richardson et N. Sime,
\emph{LEoPart: a particle library for FEniCS},
Computers \& Mathematics with Applications 81 (2021), 289--315,
\url{https://arxiv.org/abs/1912.13375}. Voir aussi la méthode
\texttt{particles.increment},
\url{https://leopart.readthedocs.io/en/stable/tree/reference.html}.
\bibitem{yoshida} H. Yoshida, \emph{Construction of higher order symplectic
integrators}, Physics Letters A 150 (1990), 262--268,
\url{https://doi.org/10.1016/0375-9601(90)90092-3}.
\bibitem{crews} D. W. Crews, \emph{Analysis of tensor-product discontinuous
Galerkin operators for Vlasov--Poisson simulations and GPU implementation
on Python}, 2022, \url{https://arxiv.org/abs/2202.13532}.
\bibitem{xue} H. Xue et A. Zanna, \emph{Explicit volume-preserving splitting
methods for polynomial divergence-free vector fields}, 2012,
\url{https://arxiv.org/abs/1205.1935}.
\end{thebibliography}
\end{document}
